% page 625 of the facsimile (image p-624 of the scan)
% block 165.1 x 250.6 mm ; left margin 36.9 mm
\pgc{15.240mm}{0.000mm}{
\l{\cel{54}617}
\l{}
\l{}
\l{vilo. (\sou{anat.}) Singla ek la pili sur la surfaco di ula membrani}
\l{\cel{3}(mukozo di la stomako, e c.) - EFIS.}
\l{}
\l{vilajo. Mikra grupo de domi di rurani, qua ne havas stradi rekta-}
\l{\cel{3}linea, nek limito cirkum-muro o ramparo-cirkuito. - EFI.}
\l{}
\l{\sou{vilanelo}. Kansono di vilajani, maxim-multa-kaze kun kupleti}
\l{\cel{3}tri-versa e refreno, finanta per quatreno. - DEFIRS.}
\l{}
\l{\sou{vildo}. Animali quin onu kaptas chase (perdriki, qualii, lepori,}
\l{\cel{3}kunikli, kapreoli). - D.}
\l{}
\l{\sou{vilejar}. (\sou{netrans}.)Sejornar en ruro dum la sezono bela.- FI.}
\l{}
\l{\sou{vilozito}. (\sou{anat.}) Singla ek la mikra asperaji o saliaji qui}
\l{\cel{3}kovras ula surfaci quin li igas aspektar piloza.(kom ex.:}
\l{\cel{3}la viloziti intestinala). - EF.}
\l{}
\l{\sou{vimplo}. I. Telo-tuko qua, en la kostumi di la regulierini, ko-}
\l{\cel{3}vras la pektoro, la kolo, e cirkondas la infra parto di la}
\l{\cel{3}vizajo. - II. Kamizeto sen maniki (ek tulo, dentelo, e c.)}
\l{\cel{3}quan +weras la mulieri kun robo dekoltita. - E.}
\l{}
\l{\sou{vino}. Suko fermentacita di vitoberi - drinkajo alkoholoza.}
\l{\cel{3}- DEFIRS.}
\l{}
\l{\sou{vinagro}. Liquoro qua saporas acide e pikante, ek la acido-}
\l{\cel{3}fermentaco de vino. - EFS.}
\l{}
\l{\sou{vindar}. (\sou{trans.})\cel{2}Elevar grand-esforce (kargajo) per mashino}
\l{\cel{3}qua konsistas ye arboro o cilindro horizontala qua rotacas}
\l{\cel{3}sur axo ed an-cirkum qua spulesas kordego an qua ligesas la}
\l{\cel{3}kargajo. - DE.}
\l{}
\l{\sou{vinkar}. (\sou{trans.}) I. Igar su superesar (a) pri forteso, sua}
\l{\cel{3}enemikon (en batalio, milito), (b) pri habileso, talento,}
\l{\cel{3}sua konkurencanton. - II. Superpasar to quo esas obstaklo.}
\l{\cel{3}- EFIS.}
\l{}
\l{\sou{vintro}. La sezono maxim kolda di la yaro, en qua la jorni esas}
\l{\cel{3}maxime kurta, qua sequas autuno e pre-iras printempo. - DE.}
\l{}
\l{\sou{vinyeto}. (\sou{tekn.)}\cel{4}Ornamento qua reprezentas vito-branchi inter-}
\l{\cel{3}plektita, supre di la pagino unesma di libro o di chapitro.}
\l{\cel{3}Omna ornamento di la frontispico o di la pagini di libro.}
\l{\cel{3}Graburo cirkondata da kartusho, Ornamento di la kovrilo di}
\l{\cel{3}libro. Ornamento cirkum naztuko. Ornamento di paperfolii}
\l{\cel{3}letrala. -\cel{2}DEFIRS.}
\l{}
\l{\sou{violo}. I. (\sou{bot.)}\cel{3}Planto ek la familio \textquotedbl{}violacei\textquotedbl{}, kun mikra}
\l{\cel{3}floro di qua la odoro esas dolca, e la koloro inter reda e}
\l{\cel{3}blua.- L.\cel{1}\sou{viola}. II. La floro di ta planto. - EFIRS.}
\l{}
\l{\sou{violacar.(trans.}) Domajar, lezar (ulo) quan onu ya devas respek-}
\l{\cel{3}tar. - Koitar\cel{5}muliero kontre elua volo.- EFIS}
}
